% Figure: load-cell characterisation case study.
% Left panel: measured output (mV/V) vs applied load with a straight-line
% fit through it; the eye sees a good straight line.
% Right panel: the residual of the straight-line fit vs load, revealing
% a systematic bow (the nonlinearity the straight line hides).
\begin{figure}[htbp]
\centering
\resizebox{\linewidth}{!}{%
\begin{tikzpicture}[scale=1.0,
  ax/.style={->, draw=engblack, thick},
  tick/.style={draw=engblack!60, thin},
  fitline/.style={draw=engblue, thick},
  pt/.style={circle, fill=engred, inner sep=1.1pt},
  zero/.style={draw=engblack!55, thin, dashed},
  ann/.style={font=\footnotesize, align=center},
]

% --- Left: output vs load with straight-line fit ---
\begin{scope}[xshift=0cm]
  \draw[ax] (0, 0) -- (6.4, 0) node[below right, ann] {load $F$ (kg)};
  \draw[ax] (0, 0) -- (0, 4.4) node[left, ann] {output (mV/V)};
  \foreach \x/\xl in {1/100, 2/200, 3/300, 4/400, 5/500} {
    \draw[tick] (\x, -0.05) -- (\x, 0.05);
    \node[ann, below, yshift=-2pt] at (\x, 0) {\xl};
  }
  \foreach \y/\yl in {1/0.8, 2/1.6, 3/2.4, 4/3.2} {
    \draw[tick] (-0.05, \y) -- (0.05, \y);
    \node[ann, left, xshift=-2pt] at (0, \y) {\yl};
  }
  % straight-line fit: output = 0.78 * F/100 (units: 1 chart-unit = 0.8 mV/V)
  \draw[fitline, domain=0:5.4, samples=2] plot (\x, {0.765*\x});
  % data: slightly bowed (parabolic) about the line
  \foreach \x/\y in {0/0.0, 0.5/0.40, 1/0.79, 1.5/1.17, 2/1.55,
    2.5/1.93, 3/2.31, 3.5/2.68, 4/3.05, 4.5/3.42, 5/3.80} {
    \node[pt] at (\x, \y) {};
  }
  \node[ann, anchor=north west, fill=white, inner sep=1pt] at (0.2, 4.2)
    {(a) output vs load,\\straight-line fit};
\end{scope}

% --- Right: residual of straight-line fit (bowed) ---
\begin{scope}[xshift=8.5cm]
  \draw[ax] (0, -1.6) -- (6.4, -1.6) node[below right, ann] {load $F$ (kg)};
  \draw[ax] (0, -2.4) -- (0, 1.6) node[left, ann]
    {residual (\% FS)};
  \foreach \x/\xl in {1/100, 2/200, 3/300, 4/400, 5/500} {
    \draw[tick] (\x, -1.65) -- (\x, -1.55);
    \node[ann, below, yshift=-2pt] at (\x, -1.6) {\xl};
  }
  \foreach \y/\yl in {-1/-0.5, 0/0, 1/0.5} {
    \draw[tick] (-0.05, \y) -- (0.05, \y);
    \node[ann, left, xshift=-2pt] at (0, \y) {\yl};
  }
  \draw[zero] (0, 0) -- (5.4, 0);
  % residual traces a downward bow: r(F) ~ 0.5 - (F-2.5)^2 * scale (sign:
  % straight line over-reads in the middle, so residual dips)
  \foreach \x/\y in {0/0.0, 0.5/-0.55, 1/-0.92, 1.5/-1.10, 2/-1.08,
    2.5/-0.95, 3/-0.78, 3.5/-0.58, 4/-0.30, 4.5/0.10, 5/0.55} {
    \node[pt] at (\x, \y) {};
  }
  \draw[fitline, domain=0:5.0, samples=40, smooth]
    plot (\x, {-0.205*(\x)*(5-\x) + 0.14*\x});
  \node[ann, anchor=north west, fill=white, inner sep=1pt] at (0.2, 1.4)
    {(b) residual reveals\\a systematic bow};
\end{scope}

\end{tikzpicture}%
}%
\caption[A load cell characterised from data]{Characterising a
strain-gauge load cell. Panel (a) plots the measured bridge output in
millivolts per volt of excitation against applied calibration load,
with a least-squares straight line through the eleven points. To the
eye the fit is excellent; a straight-line model with one slope and
one offset looks complete. Panel (b) plots the residual of that fit,
the measured output minus the fitted line, as a percentage of full
scale. The residual is not random scatter about zero; it traces a
systematic bow, the signature of a curvature the straight line
cannot represent. The bow peaks near mid-range at roughly $0.5\,\%$
of full scale, which for a precision cell is many times the
specified linearity. The residual plot is the diagnostic that
promotes the model from a one-parameter line to a two-parameter
quadratic; reading panel (a) alone would have certified a model that
is wrong where it matters.}
\label{fig:vol02:ch02:loadcell-characterisation}
\end{figure}
